\chapter{先看几个小合数：六平均为什么值得单独讲}\label{ch:small}
\section{两个新问题：用哪个模数，能缩放哪些倍数}
前篇已经解释过，把许多位置压成三种重数后，可以用二乘二矩阵记录返回。合数平均继续使用这一工具，但有两件事要重新确定。第一，总位置数的某些素因子已经允许出现在分母里，它们不该再成为障碍。第二，已经做出了加法式的剪切，还不一定做得出完成归约所需的乘法式缩放。

六平均十元清楚地展示第一件事，十平均与十二平均十七元展示第二件事。它们都已经完成，适合当作一般构造的具体模型。四平均七元、八平均十三元和九平均十三元的意义，则是提前展示后面会出现的参数族。

\section{六平均十元：允许偶数的差分}
十元本原零和输入满足 $G\mid10$。六平均允许分母二和三，所以合法的 $G$ 是一或二；五和十不合法。\src{S10}

例如 $(1^5,(-1)^5)$ 的 $G=2$。先选三份一、三份负一，产生六个零；再选剩下四个非零位置和两个零，十项便全为零。这个两步例子直观说明：排除偶数 $G$ 会排除真实可解的输入。与之相对，$(1^9,-9)$ 的 $G=10$，模五始终全为一，确实不可解。

先留四个位置，平均另外六个，得到六份 $a$。保留位置中可以安排两个模五不同的值，于是能选一个 $w\not\equiv a\pmod5$。再平均三份 $a$ 和其余三个保留值，就进入
\[
K(u,v)=(u^6,v^3,-6u-3v).
\]
总和为零给出 $u=(-3a-w)/6$、$v=a$，因此
\[u-v\equiv(a-w)/6\ne0\pmod5.
\]
取 $x=u,y=u-v$，参数清分母并本原化后，
\[G(K)=\gcd(y,10).
\]
我们要搬动的列向量满足 $\gcd(x,y)=1$ 和 $5\nmid y$。它完全可能有偶数 $y$。若强迫 $y$ 在模十下可逆，就会误删全部 $G=2$ 的合法输入。因此这里的同余群应该保留模五条件。

\section{一个两步返回补齐了正向逆}
三个一步返回分别选择 $(u^3,v^3)$、$(u^5,w)$、$(u^3,v^2,w)$，其中 $w=-6u-3v$。它们在差坐标上是
\[
A=\frac12\begin{pmatrix}2&-1\\0&-1\end{pmatrix},\quad
B=\frac16\begin{pmatrix}-4&3\\-10&9\end{pmatrix},\quad
C=\frac16\begin{pmatrix}-4&1\\-10&1\end{pmatrix}.
\]
再做一个两步返回：先平均 $u^4,v,w$，得到六份 $t=-(u+v)/3$；然后平均三份 $t$、一份 $u$、两份 $v$。留下三份 $t$ 和一份 $u$，所以
\[
R=\frac16\begin{pmatrix}1&-1\\5&-3\end{pmatrix}.
\]
矩阵计算给出
\[
R^4=-\frac1{324}I,\qquad(AR)^2=\frac1{36}I,
\qquad(RAC)^2=\frac1{216}I.
\]
这些等式的用途与前篇相同：把一个因子移到等式另一侧，就把它的射影逆写成只向前执行的返回。例如 $A^{-1}\sim RAR$。这样，$A,C,R$ 都有正向逆，进而
\[\Sigma=C^{-1}B=\begin{pmatrix}1&-1\\0&-1\end{pmatrix}
\]
也能执行。注意六重块和三重块大小不同，$\Sigma$ 的来源是这一返回词。

记 $U(z),L(z)$ 为上、下剪切，$\Delta(z)=\diag(1,z)$。因为 $\Sigma A=\Delta(1/2)$，我们有 $D=\Delta(2)$ 及其逆，从而
\[\Sigma D\Sigma D^{-1}=U(1/2).
\]
整数幂与共轭给出所有二进上剪切。接下来算
\[
F=U(-1)C=\begin{pmatrix}1&0\\-5/3&1/6\end{pmatrix},\qquad
DFD^{-1}F^{-1}=L(-5/3).
\]
将后者的逆执行三次，得到 $L(5)$。原来需要的模五条件，恰好出现在下剪切的系数里。进一步由 $L(5/3)F=\Delta(1/6)$ 得到三的缩放，最后可以执行
\[U(\Z[1/6]),\qquad L(5\Z[1/6]),\qquad
\Delta(\Z[1/6]^\times).
\]

\section{把任意合法输入送到一个看得见的终点}
前篇的陪集回路方法此时只需处理 $\mathbf P^1(\F_5)$ 的六个点。代表取 $\infty,0,\pm1,\pm2$；闭合回路的主元都是 $\pm1,\pm2$，已经可以缩放。于是整个 $\Gamma_0(5)$ 都有正向实现。

对合法本原列 $(x,y)$，求 $\alpha x+\beta y=1$，再调整系数使 $5\mid\alpha$。这是可行的，因为调整 $\alpha\mapsto\alpha+jy$ 时，$y$ 在模五下可逆。于是
\[
\begin{pmatrix}y&-x\\\alpha&\beta\end{pmatrix}
\binom{x}{y}=\binom01
\]
给出所需归约。核心变成 $(0^6,v^3,-3v)$，再选三个 $v$、一个 $-3v$ 和两个零做一次六平均，就结束了。

这证明 $H(6)=10$。九元的固定网络则证明 $M(6)=9$。两个结果回答不同的问题：九元研究如何同时混合全部输入，十元研究同余障碍是否是唯一障碍。

\section{十平均、十二平均十七元：剪切够了，缩放还不够}
两种核心是 $(u^{10},v^6,w)$ 和 $(u^{12},v^4,w)$。它们共用两步入口；因为十七是唯一危险素数，合法域都是 $17\nmid y$。\src{S11}

两步返回的物理形式也相同：平均指定位置得到 $k$ 份 $t$，保留 $r$ 份 $t$ 和一个未参加第一步的旧值 $z$，再平均其余 $k$ 个位置。由总和零，输出是
\[u'=-(rt+z)/k,\qquad v'=t,\qquad w'=z.
\]
十平均的一组周期返回为
\[
R=\frac1{50}\begin{pmatrix}16&-6\\51&-16\end{pmatrix},
\quad E=\frac1{50}\begin{pmatrix}-35&6\\-85&16\end{pmatrix},
\quad R^2=-I/50,\quad(RE)^2=I/2500.
\]
十二平均的一组为
\[
R=\frac1{36}\begin{pmatrix}2&-2\\17&-8\end{pmatrix},
\qquad R^4=-I/5184.
\]
配合其他几个同型返回，二者都得到 $\Delta(2)^{\pm1}$、$\Sigma$、$U(\Z[1/2])$ 和 $L(17\Z[1/2])$。原稿保存全部选取表及推导，附录的来源包一并保留。

然而模十七的带符号二幂只占八种非零余数，全部非零余数有十六种。要继续照搬陪集回路证明，还需要新倍率。十平均从返回比值得到 $\Delta(5)$，并先构造 $\Delta(25)^{-1}$，再用
\[\Delta(5)^{-1}=\Delta(25)^{-1}\Delta(5)
\]
得到它的逆。十二平均以同样思路补齐三的缩放。最终工作环分别是 $\Z[1/10]$ 和 $\Z[1/6]$。

现在可选模十七单位代表
\[
\{\pm1,\pm2,\pm4,\pm5,\pm8,\pm10,\pm20,\pm40\},
\]
或
\[
\{\pm1,\pm2,\pm3,\pm4,\pm6,\pm8,\pm12,\pm24\}.
\]
每套都覆盖全部十六种余数，并且代表本身是可执行的单位。前篇的回路编译与贝祖输送于是完成证明：$H(10)=H(12)=17$。固定网络另外给出 $M(10)=M(12)=16$。

\section{小例子怎样汇入一般族}
四平均七元对应下一讲的 $t=2$，九平均十三元对应 $t=3$。八平均十三元则对应后面的 $k=2s^2,s=2$。它们已经有无条件证明；一般族要解决的问题，是怎样在每个参数处都找到足够的可逆倍率。

这些例子也说明，发现程序与最终论证可以各有明确位置。程序先找到少量返回，证明随后只使用显示的矩阵恒等式、位置选取和完整群覆盖。原始运行中的千步样本说明编译器能输出实际路径，却不能说明路径最短；后面的参数族也尚未给出统一的最短词算法。
